\begin{abstract}
Modern provenance-based Intrusion Detection Systems (IDS) for detecting Advanced Persistent Threats (APTs) suffer from a deployment bottleneck: they require curated, pristine benign training data to build normal profiles. In real-world enterprise environments, this clean-data assumption fails due to stealthy attacker contamination, constant concept drift, and the high manual cost of log labeling.

We present \textbf{ZeroCausal}, a novel provenance-based IDS that achieves zero-label anomaly detection without requiring clean training data, historical labels, or offline retraining. ZeroCausal leverages \textit{causal invariances}---stable cause-effect mechanisms inherent in system execution---which are discovered \emph{online} directly from raw, unlabeled, and potentially contaminated system logs using PCMCI conditional independence testing. Deviations are quantified by a Hybrid Causal Anomaly Score (H-CAS) combining SCM residual errors and causal p-values, and approximately-guaranteed false-alarm rates are empirically controlled via online conformal prediction. We evaluate ZeroCausal on five diverse benchmarks: three core benchmarks---one using real enterprise host logs (OpTC) and two using high-fidelity simulated environments (DARPA TC3 and NODLINK) designed to model known APT scenarios---plus two additional evaluation datasets (BETH and StreamSpot), comparing against Isolation Forest, LOF, One-Class SVM, and an unsupervised autoencoder. On OpTC, ZeroCausal achieves mean AUC $\mathbf{0.727 \pm 0.071}$ over 10 independent seeds, outperforming Isolation Forest ($0.675 \pm 0.024$), while holding its empirical alarm rate to $4.1\%$ against a 5\% budget, yielding an operational recall of $24.21\%$ and F1-score of $0.3286$. ZeroCausal achieves its strongest absolute result on BETH (AUC 0.9656), demonstrating cross-platform generalization. We also report an honest failure case on StreamSpot (AUC 0.4991), where highly heterogeneous graph-type provenance violates the linear SCM assumptions. A contamination sweep (0--20\% attack events injected in training) reveals that ZeroCausal's causal \emph{discovery} backbone is stable under sparse poisoning, but its empirical-CDF calibration step degrades when contaminated windows reach the calibration set---a finding that localises the hardening target for future work. In the drift setting, ZeroCausal's strength is clear: when the normal activity rate doubles mid-stream, static Isolation Forest and Autoencoder baselines spike to 40\% and 57\% false-alarm rates respectively; ZeroCausal's change-point detector triggers an online SCM refit and conformal queue reset, holding the false-alarm rate to 9.2\%---4--6$\times$ lower than the static alternatives. A 13$\times$ latency reduction confirms runtime readiness for streaming evaluation, although high-dimensional environments require variance-based feature selection to mitigate one-time PCMCI discovery bottlenecks.
\end{abstract}

\begin{IEEEkeywords}
Intrusion Detection, Advanced Persistent Threats, Causal Discovery, Conformal Prediction, Provenance Graphs
\end{IEEEkeywords}
